% !TeX root = ../main.tex
%% =======================================================================
%  CAPITULO 9 — CONCLUSIONES

\chapter{Conclusiones}
\label{cap:conclusiones}

El receptor DP-QPSK de \SI{128}{\giga\bit\per\second} desarrollado en este
trabajo compensa el offset Doppler y su deriva, la rotaci\'on del estado de
polarizaci\'on y el ruido de los l\'aseres del enlace LEO--GEO. Frente a cada
uno de esos deterioros del canal queda margen sobre el valor f\'isico. Tambi\'en
se desarroll\'o el balance de potencias, llegando a la conclusi\'on de que es posible
con un sistema de apuntamiento adecuado.

\section{Conclusiones generales}
\label{sec:con-generales}

Los cinco objetivos espec\'ificos planteados en la \cref{sec:intro-objetivos}
se cumplieron:

\paragraph{Modelo de la cadena}
Se model\'o la cadena completa, desde el transmisor hasta la digitalizaci\'on,
pasando por el canal satelital, con los par\'ametros de los equipos del
laboratorio. Con todo el procesamiento incluido, la simulaci\'on cruza el
umbral HD-FEC a medio dB del l\'imite te\'orico AWGN.

\paragraph{Procesamiento digital}
La cadena de siete bloques del \cref{cap:dsp} compensa todas las degradaciones
que impone el enlace. El l\'imite de captura de frecuencia es fijado por
el plegado de Nyquist del tono de cuarta potencia, no por el ruido, resultado en $\Rs/4$. Los tres lazos
quedaron caracterizados, cada uno con su ancho de banda medido. El paso de
adaptaci\'on del ecualizador se eligi\'o como el mejor compromiso entre
tolerancia al estado de polarizaci\'on y ruido de gradiente.

\paragraph{Balance de potencia}
A la distancia del enlace, en el peor instante del pase, la OSNR que llega al
receptor queda ligada de forma directa a la potencia del l\'aser transmisor. El
umbral operativo, que suma los deterioros de vuelo al receptor simulado, se
alcanza con un l\'aser comercial y un amplificador de potencia trabajando por
debajo de su saturaci\'on. Pero el margen que queda al llevar ese
amplificador al m\'aximo punto de trabajo se consume por completo en el acoplamiento a fibra monomodo
y en el margen de sistema. El enlace cierra s\'olo si el apuntamiento alcanza
el orden del microradi\'an.

\paragraph{Tolerancias}
La \cref{tab:tol-final} resume el resultado. El receptor tiene margen frente a
todos los deterioros que impone el canal. El m\'as ajustado de ellos es el
offset Doppler. El resto, que son la tasa de chirp, la deriva del estado de
polarizaci\'on, el ancho de l\'inea de los l\'aseres y el error de reloj,
quedan \'ordenes de magnitud por encima de lo que el enlace pide.

El \'unico margen estrecho proviene del hardware del receptor. El skew tolerado
entre las ramas del front-end queda apenas por debajo de la especificaci\'on
del componente. Es adem\'as el t\'ermino que m\'as pesa en el presupuesto de
vuelo. Sin embargo, este deterioro es calibrable. Haciendo esta calibraci\'on en la integraci\'on
se convierte la penalidad de vuelo en algo despreciable.

Los l\'imites de offset y de chirp no cambian con la OSNR en ninguno de los
valores barridos, as\'i que son propiedades de los algoritmos y no se pueden
mejorar con potencia \'optica.

\paragraph{Validaci\'on experimental}
El mismo DSP, sin ninguna modificaci\'on, recupera la constelaci\'on
sobre capturas de laboratorio en las tres tasas medidas, incluida la de
dise\~no. El offset Doppler emulado se compensa hasta $\Rs/4$ en las tres, lo
que confirma el mecanismo de quiebre sobre se\~nales reales. La limitaci\'on del laboratorio
est\'a en el osciloscopio, que a la tasa de dise\~no no llega a dos muestras
por s\'imbolo.

Los barridos de potencia fijan la sensibilidad del conjunto en cada tasa. Del
lado del oscilador local muestran que el desempe\~no deja de mejorar a partir
de cierta potencia, como corresponde a un receptor coherente que pasa de estar
limitado por el ruido t\'ermico a estarlo por el shot del propio oscilador.
Dentro del rango enganchado la BER medida crece con el valor absoluto del
offset. La simulaci\'on reproduce esta pendiente a medida que se disminuye el
ancho de banda del front-end. Por tanto, la degradaci\'on viene del equipo utilizado en el laboratorio y no del DSP.
Con los anchos de banda de las hojas de datos, este efecto no limita el rango de Doppler del
enlace.

\section{Aportes del trabajo}
\label{sec:con-aportes}

\begin{itemize}
  \item Un simulador completo de la cadena DP-QPSK LEO--GEO, construido sobre
        los par\'ametros del hardware de laboratorio y validado contra el
        l\'imite te\'orico AWGN.
  \item Un DSP de siete bloques que funciona tanto sobre se\~nales simuladas como
        sobre capturas reales, en tres tasas de s\'imbolo.
  \item La explicaci\'on del l\'imite de captura Doppler como el rango del
        estimador de cuarta potencia, $\Rs/4$, verificada en simulaci\'on y en
        laboratorio.
  \item Un balance de potencia a la distancia del enlace, separado en l\'imite
        te\'orico, penalidad de implementaci\'on y deterioros de vuelo.
\end{itemize}

\section{Trabajo futuro}
\label{sec:con-futuro}

\begin{itemize}
  \item Si un enlace futuro necesita m\'as que $\Rs/4$ de captura, subir las
        muestras por s\'imbolo o reemplazar la
        cuarta potencia por otro algoritmo.
  \item Calibrar el skew del front end en la integraci\'on, por debajo de lo
        que pide el presupuesto de vuelo.
  \item En el laboratorio, medir la OSNR directamente y muestrear a dos
        muestras por s\'imbolo en la tasa de dise\~no, para trazar la curva de
        BER contra OSNR y compararla con la del \cref{cap:enlace}.
  \item Modelar el apuntamiento y el acoplamiento a fibra, que son los
        t\'erminos que deciden el margen del enlace.
  \item Implementar la correcci\'on de errores en bloque.
\end{itemize}
